\begin{abstract}

\noindent Dynamic analysis is a fundamental step in engineering for predicting vibrational responses, ensuring human comfort, and preventing failures. The Finite Element Method (FEM) has been one of the most widely used tools for this purpose. However, its classical formulation requires the use of elements with restricted geometry that are highly sensitive to geometric distortion. The Virtual Element Method (VEM) emerges as a generalization of the FEM, capable of handling polygonal elements of arbitrary topology, including non-convex geometries and collinear nodes. This study aims to investigate and implement the VEM formulation to solve the generalized eigenvalue problem associated with the free vibration analysis of planar structures. The technique's main advantage lies in its ability to operate on unstructured meshes composed of generic polygons, allowing for natural mesh transitions and easier adaptation to structural boundaries. The theoretical framework adopts the governing equations of elastodynamics, solved numerically using a program developed specifically for this purpose. The feasibility and accuracy of the VEM were validated through the analysis of three structural typologies of practical interest: a perforated trapezoidal plate, the cross-section of an earth dam, and a multi-story structural wall model. Its performance was assessed using error metrics compared against results obtained via the FEM and the Generalized Finite Element Method (GFEM).

\end{abstract}
